
\documentclass[11pt,a4paper]{article}

\usepackage[utf8]{inputenc}
\usepackage[T1]{fontenc}
\usepackage[french]{babel}
\usepackage{lmodern}
\usepackage{microtype}
\usepackage[a4paper,margin=1.65cm,headheight=15pt]{geometry}
\usepackage{amsmath,amssymb,mathtools}
\usepackage{array,tabularx,multicol,booktabs}
\usepackage{enumitem}
\usepackage[most]{tcolorbox}
\usepackage{xcolor}
\usepackage{tikz}
\usetikzlibrary{arrows.meta,calc,positioning}
\usepackage{pdflscape}
\usepackage{fancyhdr}
\usepackage{hyperref}
\usepackage{needspace}

\definecolor{fondpage}{RGB}{247,248,250}
\definecolor{encre}{RGB}{31,35,40}
\definecolor{bleu}{RGB}{40,91,135}
\definecolor{vert}{RGB}{55,125,92}
\definecolor{orange}{RGB}{178,105,42}
\definecolor{violet}{RGB}{101,77,145}
\definecolor{rouge}{RGB}{170,72,72}
\definecolor{gris}{RGB}{86,91,99}
\definecolor{bleufonce}{RGB}{28,55,120}
\definecolor{bleuclair}{RGB}{238,243,255}
\definecolor{grisclair}{RGB}{245,245,245}
\definecolor{tableline}{RGB}{45,45,45}
\definecolor{softgray}{RGB}{246,247,249}
\definecolor{deepblue}{RGB}{25,62,115}
\definecolor{accent}{RGB}{20,82,130}

\hypersetup{colorlinks=true,linkcolor=bleufonce,urlcolor=bleufonce,
pdftitle={Parcours de revision -- Probabilités conditionnelles -- Premiere specialite},pdfauthor={}}

\setlength{\parindent}{0pt}
\setlength{\parskip}{0.25em}
\renewcommand{\arraystretch}{1.13}
\setlength{\tabcolsep}{4pt}
\setlist[itemize]{leftmargin=1.25em,itemsep=0.15em,topsep=0.2em}
\setlist[enumerate,1]{label=\textbf{\arabic*.}, leftmargin=1.05cm, itemsep=0.35em, topsep=0.2em}
\setlist[enumerate,2]{label=\textbf{\alph*.}, leftmargin=0.85cm, itemsep=0.25em, topsep=0.15em}

\pagestyle{fancy}
\fancyhf{}
\cfoot{\thepage}

\newcommand{\R}{\mathbb{R}}
\newcommand{\N}{\mathbb{N}}
\newcommand{\sourceq}[1]{ {\scriptsize\textcolor{bleufonce}{(site Première q#1)}}}

\newcounter{question}
\newcounter{reponse}

\newenvironment{blocquestions}[1]{%
  \begin{tcolorbox}[enhanced,breakable,colback=bleuclair,colframe=bleufonce,boxrule=0.6pt,arc=2mm,left=2mm,right=2mm,top=2mm,bottom=1.5mm,title={\bfseries #1},coltitle=white,colbacktitle=bleufonce,before skip=3mm,after skip=3mm, fontupper=\large]
  \large
  \begin{enumerate}[leftmargin=*,label=\textbf{\arabic*.},itemsep=0.85em,topsep=0.5em]
  \setcounter{enumi}{\value{question}}
}{%
  \setcounter{question}{\value{enumi}}
  \end{enumerate}
  \end{tcolorbox}
}

\newenvironment{blocreponses}[1]{%
  \begin{tcolorbox}[enhanced,breakable,colback=bleuclair,colframe=bleufonce,boxrule=0.6pt,arc=2mm,left=2mm,right=2mm,top=2mm,bottom=1.5mm,title={\bfseries #1},coltitle=white,colbacktitle=bleufonce,before skip=3mm,after skip=3mm, fontupper=\large]
  \large
  \begin{enumerate}[leftmargin=*,label=\textbf{\arabic*.},itemsep=0.45em,topsep=0.5em]
  \setcounter{enumi}{\value{reponse}}
}{%
  \setcounter{reponse}{\value{enumi}}
  \end{enumerate}
  \end{tcolorbox}
}

\newtcolorbox{correctionbox}[1]{enhanced,breakable,colback=grisclair,colframe=black!55,boxrule=0.55pt,arc=2mm,left=2mm,right=2mm,top=2mm,bottom=1.5mm,title=\textbf{#1},coltitle=black,colbacktitle=black!12,before skip=2mm,after skip=2mm}

\newtcolorbox{methodebox}[1]{enhanced,breakable,colback=white,colframe=bleu!70!black,boxrule=0.55pt,arc=2mm,left=2mm,right=2mm,top=2mm,bottom=1.5mm,title=\textbf{#1},coltitle=white,colbacktitle=bleu!70!black,before skip=2mm,after skip=2mm}

\newcommand{\titreSujet}[2]{%
  \subsection*{#1}
  \addcontentsline{toc}{subsection}{#1}
  \textit{Référence / inspiration : #2}\par\medskip\hrule\medskip
}

\newcommand{\qcm}[4]{%
\begin{center}
\begin{tabularx}{0.96\linewidth}{@{}XX@{}}
\textbf{a.} #1 & \textbf{b.} #2\\[1mm]
\textbf{c.} #3 & \textbf{d.} #4
\end{tabularx}
\end{center}}

\newcommand{\arbreDeuxNiveaux}[6]{%
\begin{center}
\begin{tikzpicture}[x=1cm,y=1cm,>=Latex]
\coordinate (O) at (0,0);
\coordinate (A) at (2,1.25);
\coordinate (Ab) at (2,-1.25);
\coordinate (AB) at (4.3,1.85);
\coordinate (ABb) at (4.3,0.65);
\coordinate (AbB) at (4.3,-0.65);
\coordinate (AbBb) at (4.3,-1.85);
\draw[->,line width=.55pt] (O)--(A) node[midway,above left,scale=.86] {#1};
\draw[->,line width=.55pt] (O)--(Ab) node[midway,below left,scale=.86] {#2};
\draw[->,line width=.55pt] (A)--(AB) node[midway,above,scale=.86] {#3};
\draw[->,line width=.55pt] (A)--(ABb) node[midway,below,scale=.86] {#4};
\draw[->,line width=.55pt] (Ab)--(AbB) node[midway,above,scale=.86] {#5};
\draw[->,line width=.55pt] (Ab)--(AbBb) node[midway,below,scale=.86] {#6};
\node[font=\small\bfseries] at (A) [right=1mm] {$A$};
\node[font=\small\bfseries] at (Ab) [right=1mm] {$\overline A$};
\node[font=\small\bfseries] at (AB) [right=1mm] {$B$};
\node[font=\small\bfseries] at (ABb) [right=1mm] {$\overline B$};
\node[font=\small\bfseries] at (AbB) [right=1mm] {$B$};
\node[font=\small\bfseries] at (AbBb) [right=1mm] {$\overline B$};
\end{tikzpicture}
\end{center}}



\newcommand{\arbreSC}[6]{%
\begin{center}
\begin{tikzpicture}[x=1cm,y=1cm,>=Latex]
\coordinate (O) at (0,0);
\coordinate (S) at (2,1.25);
\coordinate (Sb) at (2,-1.25);
\coordinate (SC) at (4.3,1.85);
\coordinate (SCb) at (4.3,0.65);
\coordinate (SbC) at (4.3,-0.65);
\coordinate (SbCb) at (4.3,-1.85);
\draw[->,line width=.55pt] (O)--(S) node[midway,above left,scale=.86] {#1};
\draw[->,line width=.55pt] (O)--(Sb) node[midway,below left,scale=.86] {#2};
\draw[->,line width=.55pt] (S)--(SC) node[midway,above,scale=.86] {#3};
\draw[->,line width=.55pt] (S)--(SCb) node[midway,below,scale=.86] {#4};
\draw[->,line width=.55pt] (Sb)--(SbC) node[midway,above,scale=.86] {#5};
\draw[->,line width=.55pt] (Sb)--(SbCb) node[midway,below,scale=.86] {#6};
\node[font=\small\bfseries] at (S) [right=1mm] {$S$};
\node[font=\small\bfseries] at (Sb) [right=1mm] {$\overline S$};
\node[font=\small\bfseries] at (SC) [right=1mm] {$C$};
\node[font=\small\bfseries] at (SCb) [right=1mm] {$\overline C$};
\node[font=\small\bfseries] at (SbC) [right=1mm] {$C$};
\node[font=\small\bfseries] at (SbCb) [right=1mm] {$\overline C$};
\end{tikzpicture}
\end{center}}

\newcommand{\arbreDT}[6]{%
\begin{center}
\begin{tikzpicture}[x=1cm,y=1cm,>=Latex]
\coordinate (O) at (0,0);
\coordinate (D) at (2,1.25);
\coordinate (Db) at (2,-1.25);
\coordinate (DT) at (4.3,1.85);
\coordinate (DTb) at (4.3,0.65);
\coordinate (DbT) at (4.3,-0.65);
\coordinate (DbTb) at (4.3,-1.85);
\draw[->,line width=.55pt] (O)--(D) node[midway,above left,scale=.86] {#1};
\draw[->,line width=.55pt] (O)--(Db) node[midway,below left,scale=.86] {#2};
\draw[->,line width=.55pt] (D)--(DT) node[midway,above,scale=.86] {#3};
\draw[->,line width=.55pt] (D)--(DTb) node[midway,below,scale=.86] {#4};
\draw[->,line width=.55pt] (Db)--(DbT) node[midway,above,scale=.86] {#5};
\draw[->,line width=.55pt] (Db)--(DbTb) node[midway,below,scale=.86] {#6};
\node[font=\small\bfseries] at (D) [right=1mm] {$D$};
\node[font=\small\bfseries] at (Db) [right=1mm] {$\overline D$};
\node[font=\small\bfseries] at (DT) [right=1mm] {$T$};
\node[font=\small\bfseries] at (DTb) [right=1mm] {$\overline T$};
\node[font=\small\bfseries] at (DbT) [right=1mm] {$T$};
\node[font=\small\bfseries] at (DbTb) [right=1mm] {$\overline T$};
\end{tikzpicture}
\end{center}}



\newcommand{\arbreFicheProbas}{%
\begin{center}
\begin{tikzpicture}[x=0.65cm,y=0.48cm,>=Latex]
\coordinate (O) at (0,0);
\coordinate (A) at (2,1.45);
\coordinate (Ab) at (2,-1.45);
\coordinate (AB) at (4.1,2.15);
\coordinate (ABb) at (4.1,0.75);
\coordinate (AbB) at (4.1,-0.75);
\coordinate (AbBb) at (4.1,-2.15);
\draw[->,line width=.45pt] (O)--(A) node[midway,above left,font=\tiny] {$P(A)$};
\draw[->,line width=.45pt] (O)--(Ab) node[midway,below left,font=\tiny] {$P(\overline A)$};
\draw[->,line width=.45pt] (A)--(AB) node[midway,above,font=\tiny] {$P_A(B)$};
\draw[->,line width=.45pt] (A)--(ABb) node[midway,below,font=\tiny] {$P_A(\overline B)$};
\draw[->,line width=.45pt] (Ab)--(AbB) node[midway,above,font=\tiny] {$P_{\overline A}(B)$};
\draw[->,line width=.45pt] (Ab)--(AbBb) node[midway,below,font=\tiny] {$P_{\overline A}(\overline B)$};
\node[font=\tiny\bfseries,fill=fondpage,inner sep=0.7pt,anchor=south] at ($(A)+(0,0.18)$) {$A$};
\node[font=\tiny\bfseries,fill=fondpage,inner sep=0.7pt,anchor=north] at ($(Ab)+(0,-0.18)$) {$\overline A$};
\node[font=\tiny\bfseries] at (AB) [right=.4mm] {$B$};
\node[font=\tiny\bfseries] at (ABb) [right=.4mm] {$\overline B$};
\node[font=\tiny\bfseries] at (AbB) [right=.4mm] {$B$};
\node[font=\tiny\bfseries] at (AbBb) [right=.4mm] {$\overline B$};
\node[font=\tiny,anchor=west,text=gris] at (4.7,2.15) {$P(A\cap B)$};
\node[font=\tiny,anchor=west,text=gris] at (4.7,0.75) {$P(A\cap\overline B)$};
\node[font=\tiny,anchor=west,text=gris] at (4.7,-0.75) {$P(\overline A\cap B)$};
\node[font=\tiny,anchor=west,text=gris] at (4.7,-2.15) {$P(\overline A\cap\overline B)$};
\end{tikzpicture}
\end{center}
}

\newcommand{\tableauProbas}{%
\begin{center}
\renewcommand{\arraystretch}{1.25}
\begin{tabular}{c|c|c|c}
 & $B$ & $\overline B$ & Total\\ \hline
$A$ & $P(A\cap B)$ & $P(A\cap\overline B)$ & $P(A)$\\ \hline
$\overline A$ & $P(\overline A\cap B)$ & $P(\overline A\cap\overline B)$ & $P(\overline A)$\\ \hline
Total & $P(B)$ & $P(\overline B)$ & $1$
\end{tabular}
\end{center}}

\tcbset{enhanced,arc=2mm,outer arc=2mm,boxrule=0.42pt,left=1.15mm,right=1.15mm,top=0.75mm,bottom=0.75mm,boxsep=0.35mm,before skip=0.8mm,after skip=0.8mm,fonttitle=\bfseries\footnotesize,coltitle=encre,before upper=\raggedright\fontsize{7.65}{8.15}\selectfont}
\newtcolorbox{fichebox}[3][]{colback=white,colframe=#2!72!black,colbacktitle=#2!18,coltitle=#2!50!black,borderline west={0.9mm}{0pt}{#2!78!black},title={#3},drop fuzzy shadow=black!5,#1}
\newcommand{\tagcol}[2]{\begin{tcolorbox}[enhanced,colback=#1!11,colframe=#1!45!black,boxrule=0.42pt,arc=2mm,left=1mm,right=1mm,top=0.45mm,bottom=0.45mm,before skip=0mm,after skip=0.85mm,fontupper=\bfseries\scriptsize,colupper=#1!55!black]\centering #2\end{tcolorbox}}
\newtcolorbox{panel}[1]{colback=fondpage,colframe=fondpage,boxrule=0pt,arc=0mm,left=0mm,right=0mm,top=0mm,bottom=0mm,height=168mm,valign=top,before skip=0mm,after skip=0mm}
\newcommand{\titrepage}{\begin{tcolorbox}[colback=white,colframe=bleu!70!black,boxrule=0.65pt,arc=3mm,left=2.5mm,right=2.5mm,top=1.15mm,bottom=1.15mm,before skip=0mm,after skip=1.4mm,drop fuzzy shadow=black!10]
{\Large\bfseries Parcours de révision -- Probabilités conditionnelles}\hfill {\normalsize Première spécialité -- sans calculatrice}
\end{tcolorbox}}

\begin{document}

\begin{center}
{\LARGE\bfseries Corrigé -- Parcours de révision -- Probabilités conditionnelles}\\[1mm]
{\large Première générale -- spécialité mathématiques}\\[1mm]
{\normalsize Sans calculatrice}
\end{center}
\vspace{2mm}
\tableofcontents
\clearpage

\section*{Partie 1 -- Réponses aux questions flash}
\addcontentsline{toc}{section}{Partie 1 -- Réponses aux questions flash}
\setcounter{reponse}{0}
\begin{blocreponses}{Réponses -- Questions flash}
  \item Deux évènements incompatibles sont deux évènements qui ne peuvent pas se produire en même temps. Leur intersection est vide.
  \item $P_B(A)$ désigne la probabilité de $A$ sachant $B$.
  \item Si $P(B)\neq0$, alors $P_B(A)=\dfrac{P(A\cap B)}{P(B)}$.
  \item On utilise $P(A\cap B)=P(B)\times P_B(A)$.
  \item $A\cap B$ signifie « $A$ et $B$ ».
  \item $A\cup B$ signifie « $A$ ou $B$ ».
  \item L'évènement contraire de $A$ se note souvent $\overline A$.
  \item $P(\overline A)=1-P(A)$.
  \item On multiplie les probabilités des branches du chemin.
  \item On additionne les probabilités des chemins qui réalisent cet évènement.
  \item Cela signifie qu'ils sont incompatibles et que leur réunion est l'univers.
  \item Si $(B_i)$ est une partition, alors $P(A)=P(A\cap B_1)+\cdots+P(A\cap B_n)$.
  \item Ils sont indépendants lorsque $P(A\cap B)=P(A)P(B)$, ou encore lorsque $P_A(B)=P(B)$ si $P(A)\neq0$.
  \item Non. Deux évènements incompatibles ne sont pas forcément indépendants.
  \item Un évènement certain a pour probabilité $1$.
  \item L'évènement impossible a pour probabilité $0$.
  \item On peut conclure que les évènements $A$ et $B$ sont indépendants.
  \item On a $P(A\cap B)=0$.
  \item L'univers est l'ensemble de toutes les issues possibles.
  \item Deux évènements indépendants vérifient $P(A\cap B)=P(A)P(B)$.
  \item Une partition est une famille d'évènements incompatibles dont la réunion est l'univers.
  \item On écrit : « $A$ et $\overline A$ forment une partition de l'univers. D'après la formule des probabilités totales, on a : ... »
  \item L'indépendance ne donne pas une formule sur l'union. Si $A$ et $B$ sont indépendants, alors $P(A\cap B)=P(A)P(B)$.
  \item Il faut préciser que cette formule est vraie seulement si $A$ et $B$ sont indépendants.
  \item $P(\overline A)=1-0{,}35=0{,}65$.
  \item Comme $A$ et $B$ forment une partition de l'univers, $P(A)+P(B)=1$.
  \item Il faut annoncer la partition, citer la formule des probabilités totales, puis écrire la somme des probabilités des intersections.
  \item Il faut vérifier une égalité de probabilités, par exemple $P(A\cap B)=P(A)P(B)$.
  \item On conclut : « Comme $P_B(A)=P(A)$ et $P(B)\neq0$, les évènements $A$ et $B$ sont indépendants. »
  \item En général, $P(A\cap B)$ ne se calcule pas par une somme. La somme intervient pour une union, pas pour une intersection.
  \item $P(A\cap B)=P(A)P_A(B)=\dfrac35\times\dfrac23=\dfrac25$.
  \item $A$ et $\overline A$ forment une partition de l'univers. D'après la formule des probabilités totales :
  \[
  P(B)=0{,}4\times0{,}25+0{,}6\times0{,}5=0{,}1+0{,}3=0{,}4.
  \]
  \item $P(A)P(B)=0{,}2\times0{,}5=0{,}1=P(A\cap B)$. Les évènements sont indépendants.
  \item $P_B(A)=\dfrac{P(A\cap B)}{P(B)}=\dfrac{0{,}12}{0{,}4}=0{,}3$.
  \item Les probabilités issues d'un même noeud ont pour somme $1$, donc $p=0{,}3$.
\end{blocreponses}
\clearpage

\section*{Partie 2 -- Corrigé du QCM}
\addcontentsline{toc}{section}{Partie 2 -- Corrigé du QCM}
\begin{blocreponses}{Réponses -- QCM}
  \item Réponse \textbf{b}. $P(\overline A)=1-0{,}28=0{,}72$.
  \item Réponse \textbf{c}. $P(A\cap B)=P(A)P_A(B)=\dfrac12\times\dfrac35=\dfrac3{10}$.
  \item Réponse \textbf{b}. $P_B(A)=\dfrac{P(A\cap B)}{P(B)}=\dfrac{0{,}08}{0{,}2}=0{,}4$.
  \item Réponse \textbf{a}. $A$ et $\overline A$ sont contraires, donc $P(A)+P(\overline A)=1$.
  \item Réponse \textbf{b}. Pour deux évènements indépendants, $P(A\cap B)=P(A)P(B)$.
  \item Réponse \textbf{b}. Si $P_A(B)=P(B)$ et $P(A)\neq0$, alors $A$ et $B$ sont indépendants.
\end{blocreponses}
\clearpage


\section*{Partie 3 -- Corrigés des sujets type bac / E3C}
\addcontentsline{toc}{section}{Partie 3 -- Corrigés des sujets type bac / E3C}

\titreSujet{Sujet 1 -- Dessert et café}{inspiré de sujets E3C Première sur arbres à trois branches -- sans calculatrice}
\begin{correctionbox}{1. Partition}
Les évènements $M$, $T$ et $N$ correspondent aux trois choix possibles concernant le dessert : macarons, tarte ou aucun dessert. Aucun client ne prend plusieurs desserts.

Ils sont donc incompatibles deux à deux et leur réunion est l'univers. Ainsi, $M$, $T$ et $N$ forment une partition de l'univers.
\end{correctionbox}

\begin{correctionbox}{2. Arbre pondéré}
On a :
\[
P(M)=\frac{40}{100}=\frac25,
\qquad
P(T)=\frac{30}{100}=\frac3{10},
\qquad
P(N)=1-\frac25-\frac3{10}=\frac3{10}.
\]
De plus :
\[
P_M(C)=\frac34,
\qquad
P_T(C)=\frac13,
\qquad
P_N(C)=\frac16.
\]
On complète donc les branches contraires :
\[
P_M(\overline C)=\frac14,
\qquad
P_T(\overline C)=\frac23,
\qquad
P_N(\overline C)=\frac56.
\]
\begin{center}
\begin{tikzpicture}[x=1cm,y=1cm,>=Latex]
\coordinate (O) at (0,0);
\coordinate (M) at (2,1.7);
\coordinate (T) at (2,0);
\coordinate (N) at (2,-1.7);
\coordinate (MC) at (4.3,2.2);
\coordinate (MCb) at (4.3,1.2);
\coordinate (TC) at (4.3,.5);
\coordinate (TCb) at (4.3,-.5);
\coordinate (NC) at (4.3,-1.2);
\coordinate (NCb) at (4.3,-2.2);
\draw[->] (O)--(M) node[midway,above left] {$\frac25$};
\draw[->] (O)--(T) node[midway,above] {$\frac3{10}$};
\draw[->] (O)--(N) node[midway,below left] {$\frac3{10}$};
\draw[->] (M)--(MC) node[midway,above] {$\frac34$};
\draw[->] (M)--(MCb) node[midway,below] {$\frac14$};
\draw[->] (T)--(TC) node[midway,above] {$\frac13$};
\draw[->] (T)--(TCb) node[midway,below] {$\frac23$};
\draw[->] (N)--(NC) node[midway,above] {$\frac16$};
\draw[->] (N)--(NCb) node[midway,below] {$\frac56$};
\node[font=\bfseries,above=1mm] at (M) {$M$};
\node[font=\bfseries,above=1mm] at (T) {$T$};
\node[font=\bfseries,below=1mm] at (N) {$N$};
\node[right=1mm,font=\bfseries] at (MC) {$C$};
\node[right=1mm,font=\bfseries] at (MCb) {$\overline C$};
\node[right=1mm,font=\bfseries] at (TC) {$C$};
\node[right=1mm,font=\bfseries] at (TCb) {$\overline C$};
\node[right=1mm,font=\bfseries] at (NC) {$C$};
\node[right=1mm,font=\bfseries] at (NCb) {$\overline C$};
\end{tikzpicture}
\end{center}
\end{correctionbox}

\begin{correctionbox}{3. Sens des probabilités}
$P(T\cap C)$ est la probabilité que le client choisi ait pris une tarte et un café.

$P_C(M)$ est la probabilité que le client ait choisi des macarons sachant qu'il a pris un café.
\end{correctionbox}

\begin{correctionbox}{4. Probabilités d'intersections}
On multiplie les probabilités le long de chaque chemin :
\[
P(M\cap C)=P(M)P_M(C)=\frac25\times\frac34=\frac3{10}.
\]
\[
P(T\cap C)=P(T)P_T(C)=\frac3{10}\times\frac13=\frac1{10}.
\]
\[
P(N\cap C)=P(N)P_N(C)=\frac3{10}\times\frac16=\frac1{20}.
\]
\end{correctionbox}

\begin{correctionbox}{5. Calcul de $P(C)$}
$M$, $T$ et $N$ forment une partition de l'univers. D'après la formule des probabilités totales, on a :
\[
P(C)=P(M\cap C)+P(T\cap C)+P(N\cap C).
\]
Donc :
\[
P(C)=\frac3{10}+\frac1{10}+\frac1{20}=\frac6{20}+\frac2{20}+\frac1{20}=\frac9{20}.
\]
\end{correctionbox}

\begin{correctionbox}{6. Probabilité inverse}
On cherche $P_C(M)$ :
\[
P_C(M)=\frac{P(M\cap C)}{P(C)}=\frac{\frac3{10}}{\frac9{20}}
=\frac3{10}\times\frac{20}{9}=\frac23.
\]
La probabilité que le client ait choisi des macarons sachant qu'il a pris un café est donc $\dfrac23$.
\end{correctionbox}

\begin{correctionbox}{7. Indépendance}
On compare $P_C(M)$ et $P(M)$.
\[
P_C(M)=\frac23
\qquad\text{et}\qquad
P(M)=\frac25.
\]
Ces deux probabilités sont différentes. Les évènements $M$ et $C$ ne sont donc pas indépendants.
\end{correctionbox}

\titreSujet{Sujet 2 -- Atelier et inscription}{inspiré de sujets E3C Première avec tableau croisé -- sans calculatrice}
\begin{correctionbox}{1. Tableau croisé d'effectifs}
Il y a $80$ élèves dans l'atelier scientifique, dont $50$ inscrits à l'avance. Donc $30$ élèves de l'atelier scientifique ne se sont pas inscrits à l'avance.

Il y a $200-80=120$ élèves qui ne participent pas à l'atelier scientifique. Parmi eux, $30$ se sont inscrits à l'avance, donc $90$ ne se sont pas inscrits à l'avance.
\[
\begin{array}{c|c|c|c}
 & I & \overline I & Total\\ \hline
S & 50 & 30 & 80\\ \hline
\overline S & 30 & 90 & 120\\ \hline
Total & 80 & 120 & 200
\end{array}
\]
\end{correctionbox}

\begin{correctionbox}{2. Probabilités demandées}
On choisit un élève parmi $200$.
\[
P(S)=\frac{80}{200}=\frac25,
\qquad
P(I)=\frac{80}{200}=\frac25,
\qquad
P(S\cap I)=\frac{50}{200}=\frac14.
\]
\end{correctionbox}

\begin{correctionbox}{3. Calcul et interprétation de $P_S(I)$}
\[
P_S(I)=\frac{P(S\cap I)}{P(S)}=\frac{\frac14}{\frac25}
=\frac14\times\frac52=\frac58.
\]
Parmi les élèves qui participent à l'atelier scientifique, la proportion d'élèves inscrits à l'avance est $\dfrac58$.
\end{correctionbox}

\begin{correctionbox}{4. Calcul et interprétation de $P_I(S)$}
\[
P_I(S)=\frac{P(S\cap I)}{P(I)}=\frac{\frac14}{\frac25}
=\frac14\times\frac52=\frac58.
\]
Parmi les élèves inscrits à l'avance, la proportion d'élèves qui participent à l'atelier scientifique est $\dfrac58$.
\end{correctionbox}

\begin{correctionbox}{5. Comparaison avec $P(I)$}
On a :
\[
P_S(I)=\frac58
\qquad\text{et}\qquad
P(I)=\frac25.
\]
Comme ces deux probabilités sont différentes, on peut conjecturer que les évènements $S$ et $I$ ne sont pas indépendants.
\end{correctionbox}

\begin{correctionbox}{6. Démonstration de non-indépendance}
On compare $P(S\cap I)$ et $P(S)P(I)$.
\[
P(S)P(I)=\frac25\times\frac25=\frac4{25}.
\]
Or :
\[
P(S\cap I)=\frac14.
\]
Comme $\dfrac14\neq\dfrac4{25}$, les évènements $S$ et $I$ ne sont pas indépendants.
\end{correctionbox}

\begin{correctionbox}{7. Analyse d'une justification incorrecte}
La justification proposée n'est pas correcte : l'indépendance ne se teste pas en comparant $P_S(I)$ et $P_I(S)$.

Pour démontrer que deux évènements ne sont pas indépendants, on peut par exemple comparer $P(S\cap I)$ et $P(S)P(I)$, ou comparer $P_S(I)$ et $P(I)$ lorsque $P(S)\neq0$.
\end{correctionbox}

\titreSujet{Sujet 3 -- Test de contrôle}{inspiré de sujets E3C Première de dépistage / test -- sans calculatrice}
\begin{correctionbox}{1. Probabilités données}
On lit dans l'énoncé :
\[
P(D)=\frac1{20},
\qquad
P(\overline D)=\frac{19}{20},
\qquad
P_D(T)=\frac9{10},
\qquad
P_{\overline D}(T)=\frac1{10}.
\]
\end{correctionbox}

\begin{correctionbox}{2. Arbre pondéré}
On complète les branches contraires :
\[
P_D(\overline T)=\frac1{10},
\qquad
P_{\overline D}(\overline T)=\frac9{10}.
\]
\begin{center}
\begin{tikzpicture}[x=1cm,y=1cm,>=Latex]
\coordinate (O) at (0,0);
\coordinate (D) at (2,1.25);
\coordinate (Db) at (2,-1.25);
\coordinate (DT) at (4.3,1.85);
\coordinate (DTb) at (4.3,0.65);
\coordinate (DbT) at (4.3,-0.65);
\coordinate (DbTb) at (4.3,-1.85);
\draw[->,line width=.55pt] (O)--(D) node[midway,above left,scale=.86] {$\frac1{20}$};
\draw[->,line width=.55pt] (O)--(Db) node[midway,below left,scale=.86] {$\frac{19}{20}$};
\draw[->,line width=.55pt] (D)--(DT) node[midway,above,scale=.86] {$\frac9{10}$};
\draw[->,line width=.55pt] (D)--(DTb) node[midway,below,scale=.86] {$\frac1{10}$};
\draw[->,line width=.55pt] (Db)--(DbT) node[midway,above,scale=.86] {$\frac1{10}$};
\draw[->,line width=.55pt] (Db)--(DbTb) node[midway,below,scale=.86] {$\frac9{10}$};
\node[font=\small\bfseries] at (D) [right=1mm] {$D$};
\node[font=\small\bfseries] at (Db) [right=1mm] {$\overline D$};
\node[font=\small\bfseries] at (DT) [right=1mm] {$T$};
\node[font=\small\bfseries] at (DTb) [right=1mm] {$\overline T$};
\node[font=\small\bfseries] at (DbT) [right=1mm] {$T$};
\node[font=\small\bfseries] at (DbTb) [right=1mm] {$\overline T$};
\end{tikzpicture}
\end{center}
\end{correctionbox}

\begin{correctionbox}{3. Calcul de $P(D\cap T)$}
On multiplie les probabilités le long du chemin $D$ puis $T$ :
\[
P(D\cap T)=P(D)P_D(T)=\frac1{20}\times\frac9{10}=\frac9{200}.
\]
\end{correctionbox}

\begin{correctionbox}{4. Calcul de $P(\overline D\cap T)$}
On multiplie les probabilités le long du chemin $\overline D$ puis $T$ :
\[
P(\overline D\cap T)=P(\overline D)P_{\overline D}(T)=\frac{19}{20}\times\frac1{10}=\frac{19}{200}.
\]
\end{correctionbox}

\begin{correctionbox}{5. Calcul de $P(T)$}
$D$ et $\overline D$ forment une partition de l'univers. D'après la formule des probabilités totales, on a :
\[
P(T)=P(D\cap T)+P(\overline D\cap T).
\]
Donc :
\[
P(T)=\frac9{200}+\frac{19}{200}=\frac{28}{200}=\frac7{50}.
\]
\end{correctionbox}

\begin{correctionbox}{6. Calcul de $P_T(D)$}
On utilise la définition d'une probabilité conditionnelle :
\[
P_T(D)=\frac{P(D\cap T)}{P(T)}=\frac{\frac9{200}}{\frac7{50}}
=\frac9{200}\times\frac{50}{7}=\frac9{28}.
\]
\end{correctionbox}

\begin{correctionbox}{7. Interprétation}
Lorsqu'une pièce a un test positif, la probabilité qu'elle soit réellement défectueuse est $\dfrac9{28}$.

Cette probabilité est inférieure à $\dfrac12$, car $\dfrac9{28}<\dfrac{14}{28}$.
\end{correctionbox}

\begin{correctionbox}{8. Réponse à l'affirmation}
L'affirmation est fausse. Même si le test est positif, la pièce n'est pas « presque certainement » défectueuse.

En effet, on vient de trouver :
\[
P_T(D)=\frac9{28}.
\]
Donc, parmi les pièces dont le test est positif, moins d'une pièce sur deux est réellement défectueuse. Cela s'explique par le fait que les pièces défectueuses sont rares au départ.
\end{correctionbox}
\end{document}
